\subsection{Sector dependence and posterior smoothing of phase labels}
\label{ian:phase-label-sectors}

Status: proved. Distance three gives exact label determination on the restricted code sector (Proposition~\ref{ian:phase-label-distance-three-determination}) but does not alone give hardness; suitable shortened Hamming codes give the stated depth and an explicit weight (Proposition~\ref{ian:phase-label-hamming-sector}); these statements, the extraction theorem for the smoothed target (Theorem~\ref{ian:phase-label-posterior-extraction}) and the flatness transfers (Proposition~\ref{ian:phase-label-posterior-flat-transfer}) are unconditional. For a prescribed decomposition the posterior-smoothed equivalence is false in both directions (Propositions~\ref{ian:phase-label-posterior-hard-shallow-acceptance} and~\ref{ian:phase-label-constant-posterior-counterexample}); in the existential reading only the forward direction is refuted. The common component bounds of both counterexamples use \(s_U\), i.e.\ they hold modulo (BC) through Corollary~\ref{cor:resolution-product-alphabet-aa}; their anti-atom condition and all variance, loss and hardness computations are unconditional.

Throughout, conditional expectations involving a label refer to the joint law
\[
 \Pr(Z\in dz,X=x)=\pi(dz)\mu_z(x).
\]
Write \(\mu_H=\mu(\cdot\mid H)\), write \(\E_H\) for conditioning this joint law on \(X\in H\), and put
\[
 g(Z)=\E[Q(X)\mid Z],\qquad
 g_H(Z)=\E_H[Q(X)\mid Z],\qquad
 B_H=\Var_H g_H.
\]
The between-label share is \(B_H/\Var_H Q\), when the denominator is positive. As in Subsection~\ref{ian:hard-sector-instances}, \(m\) is even, \(n=2^m\), \(c_a(z)=(-1)^{a\cdot z}\), \(h\) is the paired quadratic, and \(P\) is the one-query acceptance function defined there on the whole cube \(\Omega_n\), so that \(P(c_a)=h(a)\). For an even \(d\), \(h_d\) denotes the paired quadratic on \(\mathbb F_2^d\).

\begin{definition}[Shallow approximation and label-only explanation]
\label{ian:phase-label-sector-shallow-definition}
For \(q:\Omega_N\to[0,1]\), say that \(q\) is
\((t,e)\)-shallowly approximable under \(\mu\) if
\[
 \ell_\mu(t,q)=\min_{D\in\mathcal D_t}\E_\mu|q-D|\le e.
\]
Thus the error is mean absolute error and the depth is worst-case coordinate-query depth.

A label-only explanation on \(H\), at thresholds \(b,u\), means
\[
 B_H\ge b,\qquad
 \E_H\Var(Q\mid Z,H)\le u.
\]
The second condition specifies how much variance remains unexplained. All counterexamples below that assert an explanation have \(u=0\).
\end{definition}

\begin{proposition}[Distance-three sectors determine the character]
\label{ian:phase-label-distance-three-determination}
Let \(S\subseteq\mathbb F_2^m\) have minimum Hamming distance at least three, and let
\[
 H=\{c_a:a\in S\}.
\]
For the simplex joint law \eqref{ian:phase-label-simplex},
\[
 A=\xi+e_\kappa,\qquad
 \xi\sim U_{\mathbb F_2^m},\qquad
 \kappa\sim U_{\{0,\ldots,m\}},
\]
the label \(\xi\) determines \(A\) after conditioning on \(c_A\in H\). Consequently, for every observable \(Q(c_A)\),
\[
 \E_H\Var(Q\mid\xi,H)=0,\qquad
 \Var_H\E_H[Q\mid\xi]=\Var_H Q.
\]
In particular the between-label share is exactly one whenever \(\Var_H Q>0\).
\end{proposition}

\begin{proof}
If two distinct points \(a,a'\in S\) belonged to \(\xi+\{0,e_1,\ldots,e_m\}\), then
\[
 a+a'=e_i+e_j
\]
would have Hamming weight at most two. This contradicts the distance assumption. Every label with positive conditional probability therefore has exactly one compatible point of \(S\). Linearity of \(S\) is not used.
\end{proof}

Distance three does not imply positive acceptance variance. For example, the one-dimensional code generated by
\((1,0,1,0,1,0)\) has distance three, while the paired quadratic \(h\) vanishes on it. Thus hardness requires an additional size or Fourier condition.

\begin{proposition}[Shortened Hamming sectors with explicit parameters]
\label{ian:phase-label-hamming-sector}
Let \(Q=P\), so that \(Q\) is defined on all of \(\Omega_n\) and
\(Q(c_a)=h(a)\). Put
\[
 r=\lceil\log_2(m+1)\rceil.
\]
Choose a rank-\(r\) binary parity-check matrix with \(m\) distinct nonzero columns, and let \(S\) be its kernel. Then
\[
 |S|=2^{m-r},\qquad d(S)\ge3.
\]
These are shortened binary Hamming codes; in particular one may take the once-shortened Hamming lengths \(m=2^r-2\).

Under \(\mu=U_{C_n}\), the set \(H=\{c_a:a\in S\}\) has exact weight
\[
 w=2^{-r},\qquad
 \frac1{2(m+1)}\le w\le\frac1{m+1}.
\]
Whenever the following depth range is nonempty,
\[
 0\le t\le \frac m2-r-4,
\]
one has
\[
 \operatorname{Hard}
 \left(Q,\mu_H;t,\frac3{16},\frac1{32}\right).
\]
In particular,
\[
 \ell_\mu(t,Q)\ge\frac{5}{32}\,2^{-r}.
\]
Thus the depth is \(m/2-\log_2(m+1)-O(1)\), the covariance error parameter is \(1/32\), and the sector weight is exactly \(2^{-r}\).

For the character mixture
\(\mu=(1-w_n)U_n+w_nU_{C_n}\), the same conditional law holds and the exact sector mass is
\[
 \mu(H)=2^{-r}\bigl[w_n+(1-w_n)n2^{-n}\bigr].
\]
\end{proposition}

\begin{proof}
Distinct nonzero parity-check columns exclude codewords of weights one and two. Rank gives the dimension and cardinality. The matrix is obtained by retaining a spanning collection of columns of the Hamming parity-check matrix, equivalently by shortening the Hamming code.

Let \(f=(-1)^h\). Its ambient Walsh coefficients all have magnitude \(2^{-m/2}\). Every Fourier coefficient of \(f|_S\), using the characters of the additive group \(S\), is a sum of \(2^r\) ambient coefficients. Hence every such coefficient has magnitude at most
\[
 b=2^{r-m/2}.
\]
In particular
\[
 \Var_{U_S}h
 =\frac{1-(\E_{U_S}f)^2}{4}
 \ge\frac{1-b^2}{4}.
\]

A physical query to \(c_A(z)\) is a linear query to \(A\). Each nonnull leaf of a depth-\(t\) tree therefore restricts \(A\) to an affine subspace of \(S\) of codimension at most \(t\). Expanding \(h|_S\) in group characters and subtracting its constant coefficient gives
\[
 \left|\E[h\mid\text{leaf}]-\E_{U_S}h\right|
 \le\frac b2(2^t-1).
\]
Averaging against any leaf outputs in \([-1,1]\) proves
\[
 \sup_{F\in\mathcal F_t}
 |\Cov_{U_S}(h,F)|
 \le \frac{2^{r-m/2}}2(2^t-1).
\]
The stated depth bound makes this at most \(1/32\), and also gives \(b\le1/16\), so the variance exceeds \(3/16\). Proposition~\ref{ian:hard-sector-forward} gives the absolute-error bound \(2^{-r}(3/16-1/32)\).

The weight formulas follow by counting codewords. Both constituents of the character mixture are uniform when restricted to \(H\).
\end{proof}

\begin{remark}[The six-dimensional calculation and its limitation]
\label{ian:phase-label-six-dimensional-check}
For \(m=6\), take the shortened Hamming code defined by
\[
 A_1+A_3+A_5=0,\qquad
 A_2+A_3+A_6=0,\qquad
 A_4+A_5+A_6=0.
\]
This is a \([6,3,3]\) code. Write
\(u=A_3\), \(v=A_5\), \(s=A_6\). Its points are
\[
 A=(u+v,u+s,u,v+s,v,s),
\]
and direct substitution gives \(h(A)=u\). Therefore
\[
 |S|=8,\qquad
 \mu(H)=\frac18,\qquad
 \Var_H Q=\frac14.
\]
Proposition~\ref{ian:phase-label-distance-three-determination} proves that all this variance is between simplex labels. These values also agree with direct enumeration.

This particular sector is not hard at depth one: querying \(c_A(e_3)\) determines \(Q\). The asymptotic hardness bound above has no content at \(m=6\): its depth range is empty for \(m\le16\) and first contains a positive depth at \(m=20\). The variance calculation is also sensitive to the placement of the code coordinates relative to the paired quadratic; the parameters \([6,3,3]\) alone do not specify it. Among the linear \([6,3,\ge3]\) codes, \(\Var_Hh\) takes each of the values \(0\), \(3/16\) and \(1/4\).

There is no contradiction with
Proposition~\ref{ian:phase-label-simplex-counterexample}. On \(C_n\), the simplex label explains only a vanishing fraction of the variance. On \(H\), conditioning removes every competing simplex vertex.

Finally, \(2^{-r}=\Theta(m^{-1})\) is polynomial weight in \(m\), or in the displayed depth. It is not a constant lower bound at fixed acceptance accuracy, and this calculation alone does not establish the polynomial-budget quantifiers of
Theorem~\ref{ian:hard-sector-aa-equivalence}.
\end{remark}

\begin{definition}[Posterior-smoothed acceptance]
\label{ian:phase-label-posterior-smoothed-definition}
For the original joint law, define
\[
 \widetilde q(x)=\E[g(Z)\mid X=x].
\]
For \(\mu(x)>0\), equivalently,
\[
 \widetilde q(x)=
 \frac{\displaystyle\int
       \mu_z(x)\left(\sum_y\mu_z(y)Q(y)\right)\pi(dz)}
      {\mu(x)}.
\]
Its value on \(\mu\)-null inputs may be chosen arbitrarily in \([0,1]\).

On a sector, the corresponding acceptance mean is \(g_H\), generally different from \(g\). Conditioning on \(X\in H\) preserves the posterior law of \(Z\) given each \(X=x\in H\), so
\[
 \E_H[g(Z)\mid X]=\widetilde q(X),
\]
but generally
\[
 \E_H[g_H(Z)\mid X]\ne\widetilde q(X).
\]
\end{definition}

The posterior-smoothed equivalence examined below asserts, for an admissible decomposition in the sense of Definition~\ref{ian:phase-label-certified-admissibility}, that \(\widetilde q\) fails shallow approximation with polynomial parameters if and only if \(Q\) has a polynomial-weight hard sector carrying a label-only explanation. In the prescribed reading the decomposition is fixed before either side is evaluated; in the existential reading each side may use any admissible decomposition of the same law.

\begin{theorem}[What extraction from the smoothed target actually proves]
\label{ian:phase-label-posterior-extraction}
Let \(0<e<1/8\), \(0<\eta\le1/4\), \(t\ge1\), and put
\[
 L=\left\lceil256e^{-7}\eta^{-3}\right\rceil t.
\]
If \(\ell_\mu(L,\widetilde q)>e\), then some set \(H\) satisfies
\[
 \mu(H)\ge\frac{\eta e^3}{16},\qquad
 \operatorname{Hard}
 \left(\widetilde q,\mu_H;
       t,\frac{e^2}{4},\frac{\eta e^2}{4}\right),
\]
and
\[
 \Var_H g(Z)\ge\Var_H\widetilde q(X)\ge\frac{e^2}{4}.
\]
The same holds with \(L\) replaced by the sharper depth
\(\lceil32e^{-5}\eta^{-2}\rceil t\). This is a hard sector for \(\widetilde q\), with variance of the original label observable \(g\). It need not be a hard sector for \(Q\).

Conversely, any set \(H\) with \(\mu(H)\ge w\) and
\(\operatorname{Hard}(\widetilde q,\mu_H;t,a,\gamma)\) gives
\[
 \ell_\mu(t,\widetilde q)\ge w(a-\gamma).
\]
\end{theorem}

\begin{proof}
Corollary~\ref{ian:hard-sector-explicit} applies because \(\mu\) is a law on a finite cube and \(\widetilde q:\Omega_N\to[0,1]\); its depths \(L_*\) and \(L_\dagger\) at accuracy \(e\) are the two depths above. That corollary requires no quantum realization, component simulation, or atom bound. It requires failure at the inflated depth, rather than merely at \(t\).

On the lifted sector, conditional expectation of \(g(Z)\) given \(X\) is \(\widetilde q(X)\). Total variance gives the asserted label-observable variance. The converse is Proposition~\ref{ian:hard-sector-forward}, applied to \(\widetilde q\).
\end{proof}

\begin{proposition}[Hard smoothing need not give a hard acceptance sector]
\label{ian:phase-label-posterior-hard-shallow-acceptance}
There are decompositions into product components, \(\mathcal C_\cup\)-admissible using \(s_U\), for which \(Q\) has depth one, while \(\widetilde q\) fails shallow approximation at arbitrarily large prescribed depths and fixed positive error. Thus the implication from hardness of \(\widetilde q\) to a hard sector for \(Q\) is false, in the prescribed and in the existential reading: \(Q\) has no hard sector at positive depth under any law on its inputs.
\end{proposition}

\begin{proof}
Let \(B\) be uniform on \(\mathbb F_2^m\), \(Y\) an independent fair bit, and \(R\) an independent uniform \(n\)-spin padding block. Take
\[
 X=(c_B,\,2Y-1,\,R),\qquad Q(X)=Y=\frac{1+X_{n+1}}2.
\]
Fix the label
\[
 Z=(B,U),\qquad
 U=Y\ \text{if }h(B)=1,\qquad U=*\ \text{if }h(B)=0.
\]
Every component is a product: the visible codeword is fixed, \(Y\) is fixed or fair, and \(R\) is uniform. Its min-entropy is at least \(n\log2\), so \eqref{ian:phase-label-anti-atom} holds, and, using \(s_U\), the components satisfy \eqref{ian:phase-label-component-bound} with the product polynomial
\eqref{ian:phase-label-product-polynomial} (Proposition~\ref{ian:phase-label-union-class}).

Here \(g(Z)\) is already a function of \(X\), and
\[
 \widetilde q(X)=g(Z)
 =\frac12+\left(Y-\frac12\right)h(B).
\]
Granting \(Y\) and the padding coins to a depth-\(L\) predictor leaves at most \(L\) linear queries to \(B\). The bent-function calculation gives
\[
 \ell_\mu(L,\widetilde q)
 \ge\frac{1-2^{L-m/2}}4.
\]
For every prescribed \(L\), sufficiently large even \(m\) therefore gives failure at every \(e<1/8\), including the inflated depth required in
Theorem~\ref{ian:phase-label-posterior-extraction}.

Nevertheless no law supported on these inputs can be hard for \(Q\) at depth at least one with \(\gamma<a\). Indeed \(F=2Q-1\) is a depth-one test and
\[
 \Cov(Q,F)=2\Var Q.
\]
This contradicts hardness whenever \(\Var Q\ge a\).
\end{proof}

\begin{proposition}[An exact label-only sector can have constant smoothing]
\label{ian:phase-label-constant-posterior-counterexample}
There is a decomposition satisfying \eqref{ian:phase-label-anti-atom} and, using \(s_U\), \eqref{ian:phase-label-component-bound}, with
\(\widetilde q\equiv1/2\), but with an exact label-only explanation on a weight-\(1/2\) hard sector. The sector hardness depth tends to infinity and its acceptance variance stays bounded below. Thus the reverse implication in the posterior-smoothed equivalence is false for a prescribed decomposition. It is not refuted in the existential reading: the same law has an admissible decomposition, using \(s_U\), for which \(\widetilde q=Q\) is hard.
\end{proposition}

\begin{proof}
Let \(m\ge10\) be even, let \(A_2=1\), and choose all other coordinates of \(A\) independently and uniformly. Add an independent uniform \(n\)-spin padding block \(R\), and put
\[
 X=(c_A,R),\qquad Q(X)=P(c_A),\qquad Z=(A_3,\ldots,A_m).
\]
Here \(P\) acts on the visible block only. Writing \(d=m-2\) and \(B=Z\), its value on the support is
\[
 Q=h(A)=A_1+h_d(B)\pmod2.
\]
Therefore
\[
 g(Z)=\frac12,\qquad \widetilde q(X)=\frac12.
\]

The components have a common simulation bound for every acceptance function, not just this \(Q\). Let \(\mathcal V\) be the class of laws \(\frac12(\delta_v+\delta_{v'})\otimes U_n\) on \(\Omega_{2n}\), with \(v\ne v'\) fixed visible strings. One physical query at a coordinate where \(v\) and \(v'\) differ identifies the selected visible string. Pinning the visible block turns a \(T\)-query acceptance function into a \(T\)-query acceptance function of the padding, as in the proof of Proposition~\ref{ian:phase-label-union-class}, and \(s_U\) simulates it on the uniform padding at error \(\delta\). Hence, using \(s_U\), \(\mathcal V\) satisfies \eqref{ian:phase-label-component-bound} with
\[
 \Psi_{\mathcal V}(T,y)=1+C_U(T+y)^{d_U}.
\]
The components of the present decomposition lie in \(\mathcal V\), since \(A_1\) selects one of two distinct visible strings. Component min-entropy is
\((n+1)\log2\), so \eqref{ian:phase-label-anti-atom} holds.

Take \(H=\{A_1=0\}\). It has weight \(1/2\), and on it
\[
 Q=h_d(B),\qquad
 \Var_H Q=\frac{1-2^{-d}}4.
\]
The label determines \(Q\) exactly. A visible query is now a linear query to \(B\); padding bits are independent coins. The bent Fourier calculation gives, for \(d\ge8\),
\[
 \operatorname{Hard}
 \left(Q,\mu_H;t,\frac3{16},\frac1{32}\right)
 \quad\text{whenever}\quad t\le d/2-4.
\]
Thus \(B_H=\Var_H Q\), with zero residual variance, while
\(\ell_\mu(t,\widetilde q)=0\) even at depth zero.

For the existential reading, use instead the visible-block label \(Z'=c_A\), with product components \(\delta_{c_A}\otimes U_n\) of min-entropy \(n\log2\); using \(s_U\), this decomposition is \(\mathcal C_\cup\)-admissible. Since \(Q\) is a function of the visible block, \(\E[Q\mid Z']=Q\) and the smoothed target equals \(Q\). The sector \(H\) gives \(\ell_\mu(t,Q)\ge\frac12(3/16-1/32)\) by Proposition~\ref{ian:hard-sector-forward}.
\end{proof}

\begin{proposition}[Valid transfer under flatness]
\label{ian:phase-label-posterior-flat-transfer}
Put
\[
 v=\E\Var(Q\mid Z),\qquad
 e_P=\E|Q-\widetilde q|.
\]
Then
\[
 e_P\le\sqrt v,\qquad
 |\ell_\mu(t,Q)-\ell_\mu(t,\widetilde q)|\le e_P.
\]
Consequently, a set \(H\) with \(\mu(H)\ge w\) and \(\operatorname{Hard}(Q,\mu_H;t,a,\gamma)\) gives
\[
 \ell_\mu(t,\widetilde q)\ge w(a-\gamma)-e_P.
\]
In particular, \(\sqrt v\le w(a-\gamma)/2\) gives a lower bound
\(w(a-\gamma)/2\).

Conversely, let \(t\ge1\), \(0<e<1/8\), \(0<\eta\le1/4\), and suppose
\[
 \ell_\mu\left(
 \left\lceil32768e^{-7}\eta^{-3}\right\rceil t,
 \widetilde q\right)>e,\qquad
 v\le\frac{\eta e^7}{2^{27}}.
\]
Then a set \(H\) satisfies, with \(a=e^2/16\),
\[
 \mu(H)\ge\frac{\eta e^3}{128},\qquad
 \operatorname{Hard}(Q,\mu_H;t,a,\eta a),
\]
\[
 \E_H\Var(Q\mid Z,H)\le\frac{a^2}{4096},\qquad
 B_H\ge a-\frac{a^2}{4096}\ge\frac{e^2}{32},
\]
and
\[
 |\Cov_H(g_H,F)|\le\eta a+\frac a{64}
 \qquad(F\in\mathcal F_t).
\]
The depth multiplier may be replaced by \(\lceil1024e^{-5}\eta^{-2}\rceil\).
\end{proposition}

\begin{proof}
Conditional Jensen gives
\[
 \E(Q-\widetilde q)^2
 \le\E(Q-g(Z))^2=v.
\]
The loss comparison is the triangle inequality. Combining it with
Proposition~\ref{ian:hard-sector-forward} proves the first transfer.

The displayed flatness assumption implies \(\sqrt v\le e/2\). Hence the failure for \(\widetilde q\) implies failure for \(Q\), at the same depth, with error \(e/2\). Apply
Corollary~\ref{ian:hard-sector-explicit} to \(Q\) with accuracy \(e/2\). Its depth multipliers are
\(256(e/2)^{-7}\eta^{-3}=32768e^{-7}\eta^{-3}\) for \(L_*\) and \(32(e/2)^{-5}\eta^{-2}=1024e^{-5}\eta^{-2}\) for \(L_\dagger\), its weight is at least
\(w_0=\eta e^3/128\), and its dispersion is \(a=e^2/16\). The assumed bound is exactly
\[
 v\le w_0a^2/4096.
\]
Proposition~\ref{ian:phase-label-flat-components} completes the proof.
\end{proof}

\begin{remark}[Status of the posterior-smoothed equivalence]
\label{ian:phase-label-posterior-equivalence-status}
For a prescribed decomposition, neither direction of the posterior-smoothed equivalence is valid. Both counterexamples satisfy the anti-atom condition and, using \(s_U\), the common component bound, and both survive arbitrary polynomial changes of depth parameters by increasing \(m\). In the existential reading, Proposition~\ref{ian:phase-label-posterior-hard-shallow-acceptance} still refutes the forward direction, since its \(Q\) has no hard sector at all; Proposition~\ref{ian:phase-label-constant-posterior-counterexample} does not refute the converse, since another admissible decomposition of the same law makes \(\widetilde q=Q\) hard.

The unconditional extraction equivalence concerns hard sectors for
\(\widetilde q\) itself, with the parameters in
Theorem~\ref{ian:phase-label-posterior-extraction}. Transferring it to \(Q\) requires additional information. Proposition~\ref{ian:phase-label-posterior-flat-transfer} supplies explicit polynomial relations under flatness. Component simulability alone supplies no such flatness.
\end{remark}
